\documentclass[11pt]{article}

% ---- Fonts & encoding ----
\usepackage[T1]{fontenc}
\usepackage[utf8]{inputenc}
\usepackage{mathptmx}

% ---- Page geometry ----
\usepackage[top=1in, bottom=1in, left=1in, right=1in]{geometry}

% ---- Packages ----
\usepackage{tikz}
\usetikzlibrary{shapes.geometric, arrows.meta, positioning, fit, backgrounds}
\usepackage{booktabs}
\usepackage{enumitem}
\usepackage{microtype}
\usepackage{setspace}
\usepackage{titlesec}
\usepackage{hyperref}
\usepackage{float}
\usepackage{indentfirst}

% ---- Section formatting ----
\titleformat{\section}{\normalfont\large\bfseries}{\thesection.}{0.5em}{}
\titleformat{\subsection}{\normalfont\normalsize\bfseries}{\thesubsection}{0.5em}{}
\titlespacing*{\section}{0pt}{10pt}{4pt}
\titlespacing*{\subsection}{0pt}{8pt}{2pt}

% ---- Spacing ----
\setlength{\parindent}{1em}
\setlength{\parskip}{3pt}

% ---- Header info ----
\title{\textbf{Offloaded Spatial Audio Rendering for XR:\\
       Latency, Power, and Quality Tradeoffs}\\[6pt]
       \normalsize CS 534 IC -- Progress Report -- Spring 2026}
\author{Abhi Ramakrishnan \quad Sethu Eapen\\
        \small University of Illinois at Urbana-Champaign}
\date{}

% -----------------------------------------------------------------------
\begin{document}
\maketitle
\thispagestyle{plain}

% -----------------------------------------------------------------------
\section{Introduction}
% -----------------------------------------------------------------------

Spatial audio is an important part of XR immersion. When sound does not
match the physical space a user is in, it is immediately noticeable.
Environment-aware spatial audio addresses this by adapting to the actual
geometry and surface materials of the room, producing a Room Impulse
Response (RIR) that captures how sound propagates in that specific space.

Running this pipeline on a headset is expensive. SAMOSA (Xu et al., UIST 2025),
a recent system for on-device scene-aware spatial audio, achieves 58\,ms
to compute an updated Room Impulse Response (RIR) on a Snapdragon XR2+ Gen 2
chip. This update latency matters because it determines how quickly the
acoustic model can adapt when a user moves to a new position in a space.
SAMOSA was only evaluated on a single chip, leaving open questions about
how it performs on other hardware.

A natural response to on-device cost is offloading computation to a nearby
server, which is the approach taken by RemoteVIO (Jiang et al., MMSys 2025)
for visual-inertial odometry in ILLIXR. For spatial audio, however, offloading
introduces a different constraint: audio-visual synchronization. Prior work
(Yeregui et al., 2024) measured audio latencies of 185\,ms over Ethernet
when fully offloading the XR rendering pipeline, exceeding the 160\,ms
threshold above which audio-video desync becomes noticeable to users. Full
offloading does not appear viable for audio, but a partial offload that keeps
lightweight steps on the headset and sends only the heavy ML inference to the
server might stay within budget.

We study the latency, power, and audio quality tradeoffs of different
offloading split points for a scene-aware spatial audio pipeline, using
a MacBook as the device proxy and a Jetson Orin AGX as the server.

% -----------------------------------------------------------------------
\section{Related Work}
% -----------------------------------------------------------------------

\subsection{SAMOSA}

Xu et al.\ present an on-device system for scene-aware spatial
audio in XR. The pipeline runs four stages: shoebox room geometry estimation
from depth data, pixel-wise surface material classification using DeepLabv3+,
room type classification using MobileNetV2 to further refine acoustic
parameters, and RIR synthesis via the Image Source Method. It achieves
58\,ms end-to-end with under 3\% CPU usage on a Snapdragon XR2+ Gen 2.
Our pipeline follows the same structure but omits the MobileNetV2 room
classification step, focusing on the geometry and material inference stages
and how they partition across device and server. Notably, SAMOSA reports
only aggregate latency and does not break down timing per stage; our
profiling infrastructure provides this breakdown.

\subsection{RemoteVIO}

Jiang et al.\ study offloading of head tracking (VIO) in an
end-to-end ILLIXR deployment, showing that careful partitioning of a
latency-sensitive XR workload can reduce headset power without unacceptable
quality loss. We follow the same experimental structure
applied to the spatial audio pipeline instead of VIO.

\subsection{Edge Rendering for XR}

Yeregui et al.\ fully offload XR rendering and audio to an edge
server, measuring 185\,ms audio latency over Ethernet and up to 324\,ms
over WiFi. They establish 160\,ms as the perceptual threshold above which
audio-visual desync becomes noticeable. Our work asks whether a smarter
partial offload can stay within that budget where full offloading cannot.

% -----------------------------------------------------------------------
\section{Approach and Evaluation Plan}
% -----------------------------------------------------------------------

\subsection{System Architecture}

We build a scene-aware spatial audio pipeline using open-source components.
A MacBook serves as the device proxy and a Jetson Orin AGX serves as the
server, connected over a local network. The pipeline is implemented in Python,
with each stage as a standalone module. On the device side, shoebox room
geometry is estimated from a PLY mesh or depth frame and binaural rendering
is handled via Steam Audio, which requires continuous access to the current
head pose. On the server side, a Python process (\texttt{server.py}) accepts
RGB frames, runs material segmentation, computes the RIR using
\texttt{pyroomacoustics}, and returns the result over a local network socket.

We run segmentation in two modes, selectable via a flag. In the first mode,
SegFormer-b0 (Xie et al., NeurIPS 2021) pretrained on ADE20K runs inference
and its output feeds directly into RIR computation. In the second mode,
DeepLabv3+ with a MobileNetV3 backbone (torchvision pretrained) runs inference
for timing purposes only; its output is discarded and materials are set from
a fixed default, since this model is trained on PASCAL VOC and does not produce
valid indoor surface labels. The second mode lets us measure what segmentation
latency would look like with a lighter, SAMOSA-class model on our hardware
without needing to retrain on ADE20K. Audio quality metrics are only reported
for the SegFormer mode.

\subsection{Offloading Configurations}

\begin{table}[H]
\small\centering
\caption{Offloading configurations.}
\label{tab:configs}
\begin{tabular}{@{}llll@{}}
\toprule
\textbf{Config} & \textbf{On Device} & \textbf{On Server} & \textbf{Expected tradeoff}\\
\midrule
A (baseline) & Geometry + Materials + RIR & Nothing                & Highest power, lowest latency\\
B (partial)  & Geometry + RIR             & Materials              & Moderate savings\\
C (heavy)    & Geometry                   & Materials + RIR        & Large savings, latency risk\\
D (full)     & Nothing                    & Geometry + Materials + RIR & Lowest power, likely too slow\\
\bottomrule
\end{tabular}
\end{table}

Binaural rendering via Steam Audio runs on device in all configurations, as
it requires continuous access to the current head pose and the local audio
source. Config A reproduces the SAMOSA on-device baseline. Configs B and C
are partial offloads at different granularities. Config D sends all scene
processing to the server and corresponds to the full-offload case studied
by Yeregui et al., which we expect to exceed the perceptual budget.

\subsection{Metrics}

For each configuration we measure:

\begin{itemize}[noitemsep]
  \item \textbf{End-to-end latency}: time from camera frame capture to
        rendered audio output, compared against the 160\,ms threshold.
  \item \textbf{Per-step latency}: breakdown across depth estimation,
        material inference, RIR computation, and rendering.
  \item \textbf{Network bandwidth}: bytes transmitted per frame for each config.
  \item \textbf{Audio quality}: RIR similarity against a ground truth RIR
        computed by running \texttt{pyroomacoustics} directly on ScanNet's
        known room geometry, using reverberation time ($T_{60}$) and
        normalized echo density as metrics.
\end{itemize}

Experiments run on a MacBook and Jetson Orin AGX connected over a local
network, with additional WiFi runs to test sensitivity to network conditions.

% -----------------------------------------------------------------------
\section{Methodology}
% -----------------------------------------------------------------------

The pipeline is implemented in Python with five standalone modules:
\texttt{shoebox.py} (PLY loading, PCA alignment, AABB room estimation),
\texttt{segmentation.py} (SegFormer-b0 via HuggingFace Transformers, ADE20K class mapping),
\texttt{scene\_classifier.py} (rule-based acoustic preset selection),
\texttt{acoustics.py} (pyroomacoustics ShoeBox RIR computation), and
\texttt{render.py} (FFT convolution via scipy). A \texttt{profiler.py} module
wraps every stage with \texttt{perf\_counter} timers to produce per-stage latency
breakdowns. The full pipeline is orchestrated by \texttt{pipeline.py}, which runs
shoebox estimation and material segmentation in parallel via \texttt{ThreadPoolExecutor}.

SegFormer model weights are cached at the module level so they are loaded from
disk only once per server process, reflecting the deployment scenario where the
server handles repeated RIR update requests without restarting.

Currently, the pipeline takes a pre-made PLY mesh and a static RGB render as
inputs, with source and listener positions set to fixed offsets within the room.
The next step is to replace these with inputs from ScanNet, a dataset of
RGB-D indoor scans where each scene includes per-frame depth, RGB, and camera
poses alongside a ground-truth PLY mesh. With ScanNet, the depth frame drives
shoebox estimation directly and the camera pose provides the real listener
position, removing the hardcoded placement. The ground-truth PLY mesh serves
as a reference: running pyroomacoustics on it produces a reference RIR that
we compare against our pipeline's output using $T_{60}$ and normalized echo
density as audio quality metrics.

% -----------------------------------------------------------------------
\section{Preliminary Results}
% -----------------------------------------------------------------------

Table~\ref{tab:results} shows per-stage  latency for Config A on a desktop workstation with a NVIDIA RTX 4090 (24 GB VRAM). Material segmentation dominates at 139--273\,ms,
accounting for over 95\% of total pipeline time. All other stages are under
25\,ms combined.

\begin{table}[H]
\small\centering
\caption{Per-stage latency, Config A (all on-device, CPU, warm runs).}
\label{tab:results}
\begin{tabular}{@{}lr@{}}
\toprule
\textbf{Stage} & \textbf{Latency (ms)}\\
\midrule
Shoebox estimation     & 8\,--\,16\\
Material segmentation  & 139\,--\,273\\
Scene classification   & $<$1\\
RIR synthesis          & 3\,--\,4\\
Audio rendering        & 5\,--\,6\\
\midrule
\textbf{Total}         & \textbf{155\,--\,300}\\
\bottomrule
\end{tabular}
\end{table}

The total pipeline latency of 155--300\,ms is substantially higher than SAMOSA's
reported 58\,ms end-to-end. The gap is driven almost entirely by segmentation:
our SegFormer inference takes 139--273\,ms per frame, whereas SAMOSA's fine-tuned
DeepLabv3+ MobileNetV2 is reported to run in approximately 10\,ms on dedicated
mobile hardware. The remaining stages---shoebox estimation (8--16\,ms), RIR
synthesis (3--4\,ms), and audio rendering (5--6\,ms)---are broadly comparable to
what SAMOSA would require for equivalent geometry and synthesis steps. Replacing
or offloading the segmentation module is therefore the single most impactful
change available.

\subsection{SAMOSA Emulation Mode}

To measure what the pipeline would look like with a SAMOSA-class segmentation
model, we introduced a SAMOSA emulation mode (\texttt{--samosa-mode}). In this
mode, material labels are loaded from a pre-computed SegFormer cache
(\texttt{segmentation\_cache.json}) and DeepLabv3+ with a MobileNetV3-Large
backbone (torchvision pretrained) is run solely to consume equivalent compute
time; its output is discarded. This matches the architectural profile of
SAMOSA's segmentation slot without requiring a fine-tuned model.

\begin{table}[H]
\small\centering
\caption{Per-stage latency, Config A SAMOSA emulation mode (warm pipeline runs, RTX 4090).}
\label{tab:samosa_results}
\begin{tabular}{@{}lr@{}}
\toprule
\textbf{Stage} & \textbf{Latency (ms)}\\
\midrule
Shoebox estimation                  & 6.9\,--\,9.6\\
Segmentation (MobileNetV3, SAMOSA)  & 39.7\,--\,59.7\\
Scene classification                & $<$1\\
RIR synthesis                       & 3.3\,--\,9.4\\
Audio rendering                     & 1.7\,--\,4.6\\
\midrule
\textbf{Total}                      & \textbf{53.1\,--\,82.6}\\
\bottomrule
\end{tabular}
\end{table}

The SAMOSA emulation mode total of 53--83\,ms is closely comparable to SAMOSA's
reported 58\,ms, confirming that the pipeline architecture and non-segmentation
stages are well-matched to SAMOSA's performance envelope. The segmentation slot
(39.7--59.7\,ms in these runs) is elevated relative to the true inference cost
because each \texttt{pipeline.py} invocation is a fresh process that must load
the DeepLabv3+ weights from disk before the first timed call. When the
segmentation module is run in isolation within a persistent process---as it
would be in a deployed server that handles repeated RIR update
requests---inference time drops to approximately 8.6\,ms (median, RTX 4090,
CUDA, 5 warm runs). A persistent server process would therefore bring the
segmentation slot to under 10\,ms and the total pipeline latency to under
30\,ms, comfortably within SAMOSA's 58\,ms reference and well below the
160\,ms perceptual threshold.

% -----------------------------------------------------------------------
\section{Changes from Original Proposal}
% -----------------------------------------------------------------------

\textbf{Testbed: MacBook + Jetson Orin without ILLIXR.}
The original proposal targeted ILLIXR as the device-side testbed. On instructor
feedback, full ILLIXR integration was deemed overly ambitious for the available
timeline. We instead run the pipeline directly on a MacBook (device proxy) and
Jetson Orin AGX (server) without the ILLIXR layer. This keeps the hardware
setup close to the proposal while removing the integration overhead. Power
measurements via \texttt{tegrastats} on the Jetson remain feasible with this setup.

\textbf{Segmentation model: SegFormer-b0 instead of DeepLabv3+.}
The proposal specified DeepLabv3+ following SAMOSA. SAMOSA's model was
custom-trained on COCO-Stuff and is not publicly available, and no pretrained
DeepLabv3+ trained on ADE20K exists off the shelf. We use SegFormer-b0
pretrained on ADE20K, which provides equivalent pixel-level indoor surface
classification (150 classes) with a comparable parameter count (3.75M).
Since the segmentation model is offloaded to the server in Configs B--D,
its specific architecture does not affect the offloading comparison.

% -----------------------------------------------------------------------
\section{Remaining Timeline}
% -----------------------------------------------------------------------

\begin{table}[H]
\small\centering
\caption{Remaining project milestones.}
\label{tab:timeline}
\begin{tabular}{@{}ll@{}}
\toprule
\textbf{Task} & \textbf{Target}\\
\midrule
\texttt{server.py} socket layer; device/server split working end-to-end & Apr 10\\
Config B implemented and measured (offload segmentation)                & Apr 13\\
Configs C and D implemented and measured                                & Apr 16\\
ScanNet ground truth RIR comparison; audio quality metrics              & Apr 19\\
Full experiment runs across all configs; latency and bandwidth tables   & Apr 21\\
Results analysis; bottleneck identification; plots and tables           & Apr 24\\
Final report and presentation                                           & Apr 28\\
\bottomrule
\end{tabular}
\end{table}

\end{document}
